\documentclass[11pt,a4paper]{article}
\usepackage[T1]{fontenc}
\usepackage[utf8]{inputenc}
\usepackage{lmodern}
\usepackage[margin=25mm]{geometry}
\usepackage{amsmath,amssymb,stmaryrd,booktabs,tabularx,xurl,xcolor,hyperref}
\hypersetup{colorlinks=true,linkcolor=blue!50!black,urlcolor=blue!50!black,
  pdftitle={Typed Symbolic Regions and Resumable CTL},pdfauthor={Mettapedia Working Notes}}
\setlength{\emergencystretch}{3em}
\setlength{\parindent}{0pt}
\setlength{\parskip}{5pt}
\newcommand{\code}[1]{\nolinkurl{#1}}
\title{Typed Symbolic Regions and Resumable CTL}
\author{Mettapedia Working Notes}
\date{1 November 2026}
\begin{document}
\maketitle
\begin{abstract}
Ten Lean modules connect guarded finite-state programs to symbolic Boolean
operations, independent temporal path semantics, a resumable fixed-point
solver and a shared decision-table evaluator. The forward-image and universal
predecessor operators satisfy an adjunction: a requested output predicate
determines the largest domain-qualified input set that guarantees output
safety. This is an executable, restricted refinement-inference mechanism,
not a classifier for arbitrary programs. A separate private-fold theorem
preserves ordinary results through late refusal while keeping interruption
distinct. The results expose the information that an implementation must
retain: a shared source witness, valid-state domain, fixed transition meaning,
approximation polarity and original fallback state.
\end{abstract}

\section{Scope and representation}

The symbolic state carrier is $\Sigma_n=\mathrm{Fin}(n)\to\mathrm{Bool}$.
A diagram is an existing \code{BDDCore.BDD n} paired with its ordered and
reduced invariant. Boolean operations reuse the existing BDD implementation.
No second Boolean decision-tree datatype is introduced for the logical layer.

A source rule consists of a pure Boolean guard and one Boolean expression for
each next-state bit. All expressions read the same old state. A list of rules
denotes a nondeterministic relation: a transition exists when some enabled
rule produces its target state. A separate list observer retains repeated
rule occurrences; the relational observer intentionally forgets them.

The valid encoded states form an explicit diagram $D$. The restricted relation
$R_D$ requires both endpoints to lie in $D$. The executable rule checker tests
whether
\[
 D(s)\land g(s)\land\neg D(\operatorname{next}(s))
\]
is unsatisfiable. Its equivalence with semantic invariant preservation is
proved. Successful inference returns a program carrying a proof that every
rule preserves $D$. Consequently, restricting to $D$ drops no transition
from a valid source state. Without that certificate, a domain restriction
could conceal an escaping transition rather than justify an optimization.

This is a source language of guarded bit-vector assignments, not yet a
decoder for arbitrary MeTTa atoms or the CTL benchmark's state identifiers.
The relevant executable and proof interfaces are
\code{StateTransducer.infer}, \code{checkRule_iff} and
\code{valid_transition_not_dropped}.

\section{Substitution and correlated transition images}

The substitution compiler replaces each diagram variable with a Boolean
diagram. Its theorem is simultaneous substitution:
\[
 \llbracket\operatorname{substitute}(\sigma,t)\rrbracket(s)
 =\llbracket t\rrbracket
       (v\mapsto\llbracket\sigma(v)\rrbracket(s)).
\]
This permits a symbolic predecessor to compute the effect of an entire
assignment vector without sequentially mutating its inputs. The swap control
distinguishes this operation from in-place assignment.

For a forward image, the transition graph uses an output-bit block followed
by an input-bit block. Coordinate constraints are conjoined before the input
block is existentially projected away. Thus every output coordinate shares
\emph{one} source witness. For a rule $(x,y)\mapsto(x,x)$, the reachable
outputs satisfy equality of the two coordinates. Projecting the coordinates
independently would admit incompatible outputs such as $(1,0)$; a checked
negative control excludes that mistake.

For diagrams $A,B$, the symbolic operations implement
\begin{align*}
 \operatorname{post}_{R_D}(A)
   &=\{t\mid\exists s\in A,\ R_D(s,t)\},\\
 \operatorname{must}_{R_D}(B)
   &=\{s\in D\mid\forall t,\ R_D(s,t)\Rightarrow t\in B\}.
\end{align*}
The forward operator uses the existing existential-projection machinery.
The backward operator is a domain-relative complement of an existential
predecessor. Global negation is insufficient when not every bit pattern is
a valid encoded state.

\section{Output safety as refinement inference}

\code{RelationalImage.image_must_adjunction} establishes
\[
 \operatorname{post}_{R_D}(A)\subseteq B
 \quad\Longleftrightarrow\quad
 A\cap D\subseteq\operatorname{must}_{R_D}(B).
\]
The right adjoint is therefore the largest domain-qualified set of inputs
whose successors satisfy the requested output refinement. This moves beyond
recognizing one syntactic spelling: different rule presentations with the
same transition relation induce the same semantic safety condition.

The scope of ``largest'' matters. It is exact for this finite-state relation
and this output observation, not globally optimal code, minimum description
length or a cheapest storage representation. The construction can produce
large diagrams. A profitability decision needs a separate cost model.

Output safety does not imply progress. A state with no successor satisfies
every universal output condition vacuously, including the false condition.
The negative control computes that case and contrasts it with a false
existential-successor test. A profile requiring a result must separately
require existence; an infinite-run requirement is stronger still.

The information boundary is precise: existential images and CTL state sets
may forget occurrence multiplicity, but a bag-valued source observer may not.
Likewise, source correlations, valid-state membership and transition meaning
remain necessary even when the physical representation becomes smaller.

\section{Independent temporal semantics}

The temporal reference meaning is not defined by the symbolic solver.
\code{CTLFixedPoints.Until} is an inductive finite path: a state either
already meets the goal, or satisfies the holding condition and takes a step
to another such path. \code{Globally} is an infinite sequence of states with
the holding predicate and transition relation satisfied at every index.
The formalization proves
\begin{align*}
 E[H\ U\ G]&=\mu X.\ G\cup(H\cap\operatorname{pre}_R(X)),\\
 EG\,H&=\nu X.\ H\cap\operatorname{pre}_R(X).
\end{align*}
The reverse greatest-fixed-point direction constructs an infinite path by
choosing successors inside the greatest fixed point. Its use of classical
choice is visible in the foundation audit.

Deadlocks are explicit. A deadlocked state cannot satisfy the infinite-path
meaning of $EG\,H$, even when $H$ holds there. Adding a self-loop changes that
answer. The typed transducer supplies an explicit stuttering completion and
proves that every valid state then has a valid successor. The finite-until
reference also permits the goal to hold immediately without a successor.
Connecting to a total-Kripke CTL profile must therefore select its deadlock
policy rather than silently changing the source relation.

The symbolic implementations of predecessor, until and globally refine these
independent meanings. Nested temporal-formula parsing, source-state codecs,
fairness constraints and a whole-formula compilation pipeline remain separate
interfaces. The classic symbolic model-checking construction already uses
Boolean representations and least/greatest fixed points \cite{burch}; the
contribution here is its checked integration with typed admission, resumable
observations and the repository's existing carriers.

\section{Resumption and the information in a partial answer}

A symbolic operator carries a local refinement to an independent monotone
set operator. The solver starts at bottom for a least fixed point and at top
for a greatest fixed point. Each application increments a pass counter and
uses a complete diagram-equivalence test. Equivalence is implemented by XOR
and a checked zero test; no constant-time pointer-equality claim is made.

The result type distinguishes \code{paused frame} from \code{done frame}.
For a fixed operator, the page law is exact:
\[
 \operatorname{resume}(b,\operatorname{advance}(a,f))
 =\operatorname{advance}(a+b,f).
\]
It preserves the diagram, completion status and solver pass count, including
the final stability check. These passes are not MeTTa evaluator instructions
or the source benchmark's round counter.

Completion implies a fixed point. Bottom and top invariants identify the
completed value as the least and greatest fixed point respectively. Finite
state spaces provide ascending and descending chain conditions, so there
exists sufficient fuel for completion. This is an eventual-completion theorem,
not a numeric bound on runtime or on the cost of one Boolean operation.

Partial answers have opposite information polarity:
\begin{center}
\begin{tabularx}{\linewidth}{@{}lXX@{}}
\toprule
Iteration & Already certified & Still unresolved before completion\\
\midrule
Least fixed point & States already included & States not yet included\\
Greatest fixed point & States already excluded & States still retained\\
\bottomrule
\end{tabularx}
\end{center}
The two observation functions return \code{Option Bool} and prove soundness
of every returned Boolean. At zero fuel, an empty least approximation and a
full greatest approximation both yield unresolved observations where their
raw bits alone would be misleading.

The operator and its domain remain fixed across pages. A control shows that
changing the relation can invalidate a completed answer. A native task must
therefore retain a captured relation or qualify reuse against its revision;
the page law alone is not a concurrent-mutation theorem.

\section{A shared decision table with independent evaluation}

\code{DecisionTable} stores constants and tests whose two child indices refer
to earlier rows. Its direct evaluator caches one Boolean per row. An
independent decoder expands the records into existing BDD trees; the main
theorem proves equality between direct evaluation and the decoded tree
semantics. Invalid references return \code{none}, not false, including an
invalid child on a branch that the current input would not select.

Append-only evaluation preserves old indexed values. The counted evaluator
visits exactly one record per row on successful evaluation. A nontrivial
diamond contains six records but unfolds to eleven tree nodes. This is a
record-visit result, not a linear CPU-time result for the list-based reference
implementation: append and index lookup costs are not erased. Nor does
diagram node count measure represented state count; a one-record true
diagram represents all eight assignments to three bits.

The table is a real shared carrier, but the current package does not yet
compile arbitrary BDD trees into hash-consed tables, run memoized Boolean
operations over those tables, or verify native collection and relocation.
Those are the next representation refinements.

\section{Private fold fallback and interruption}

\code{Machines.PrivateScalarFold} reuses the existing source fold carrier.
A native step can return a private next accumulator or refuse. The local
refinement requires every accepted step to equal one pure, single-answer
source callback at the captured world. Induction then proves equality for
an entirely accepted fold, including its unchanged world.

If a later step refuses, the dispatcher runs the ordinary source fold from
the original accumulator and world, over the original ready items. It does
not continue from a partially updated native accumulator. An interruption
has a separate outcome and never initiates ordinary replay. The controls
include negative arithmetic, late signed-range overflow, cancellation with
an effectful fallback and a refuted purity obligation.

The signed-range control uses exact integers with an explicit result bound
$[-2^{63},2^{63})$; it is not a proof of an Int64 ABI or of the runtime's
operand-admission checks. The general theorem starts after item and initializer
production. It does not permit replaying their effects. Actual signed/symbol
operations, conditional and tuple extraction, callability checks, program
revisions, arenas and publication still need source and native bridges.

For lifetime design, this identifies a concrete obligation: the original
fallback inputs must remain available until the attempt's disposition is
settled. Private progress does not make them unreachable merely because it
has produced a newer accumulator.

\section{Checked modules and remaining interfaces}

The ten mathematical modules are:
\begin{itemize}
\item \code{Algebra.Order.FixedPointIteration} and \code{Logic.CTLFixedPoints};
\item \code{Logic.BDD.StatePredicates}, \code{StateTransducer},
  \code{RelationalImage}, \code{FixedPointSolver}, \code{SymbolicCTL},
  \code{DecisionTable} and \code{SymbolicCTLControls};
\item \code{Machines.PrivateScalarFold}.
\end{itemize}
\code{Logic.BDD.SymbolicRegionAudit} imports the package, prints the main
theorems' axiom dependencies and emits declaration-level foundation manifests.
The audit checks 563 declaration bodies, with ten nonlogical recursion
companions reported separately. There are no admissions or added axioms; transitive
primitive dependencies are confined to \code{propext},
\code{Classical.choice} and \code{Quot.sound}. Strict file checks treat warnings
as errors. This is a targeted package audit, not a whole-repository audit.

The immediate source task remains certified extraction of the actual scalar
fold, including its signed arithmetic, inert symbols, branches and flat
tuples. The symbolic task needs source codecs and formula compilation before
it can stand in for a source CTL computation. Native layouts, revision checks,
allocation failure and measured plan choice then refine the two routes
without changing their observation contracts.

\begin{thebibliography}{9}
\bibitem{burch}
J. R. Burch, E. M. Clarke, K. L. McMillan, D. L. Dill and L. J. Hwang.
\emph{Symbolic Model Checking: $10^{20}$ States and Beyond}.
Information and Computation 98(2), 142--170, 1992.
\href{https://doi.org/10.1016/0890-5401(92)90017-A}{doi:10.1016/0890-5401(92)90017-A}.
\href{https://mcmil.net/pubs/IC92.pdf}{Author-hosted full text}.
\end{thebibliography}
\end{document}
